\documentclass[10pt]{article}
\usepackage[utf8]{inputenc}
\usepackage{amsmath, amssymb}
\usepackage{float}
\usepackage{hyperref}
\usepackage{booktabs}
\usepackage{graphicx}
\usepackage{algorithm}
\usepackage{algpseudocode}

\usepackage[
backend=biber,
bibstyle=authoryear,
citestyle=authoryear, 
sorting=anyt,
maxcitenames=2
]{biblatex}
\addbibresource{refs.bib}

\title{Evaluating the prediction of the toxicity of plastic chemicals}
\author{Amelie Hübner}
\date{\today}

\begin{document}
\maketitle

\section{Introduction}



\section{Dataset and problem setting}
% Problem setting
[plastik ist böse, hinführung zu degradation und additives sind böse, polyurethan domaine recycling ist schwer (chemisch)]
The impact and scale of plastic production and consumption on the environment is steadily increasing. 
In 2020, cumulative primary plastic production surpassed 10,000 million tonnes. 
Suprisingly, only around 10\% of global plastic waste has ever been recycled.
Of the remaining 90\%, 14\% have been incinerated, producing greenhouse gas emissions and releasing gaseous pollutants in the process.
The remaining 76\% remain intact, leaching toxic chemicals and microplastics into the environment \parencite{walker}.


The dataset was released by \cite{monclus_mapping_2025} and contains 16.325 known plastic chemicals.
% Datensatz wird gut beschrieben, Relevanz des Themas wird deutlich, mathematisch sauber

\section{XGBoost}
% Modellannahme, Parameterschätzung und Inferenz mathematisch sauber erklärt.
% Eventuelle Vorverarbeitung der Daten und mögliche Optimierungen erklärt.
% Grenzen des Modells aufgezeigt. Pseudocode für Lernverfahren.

% Modellstruktur und Modellannahmen

% Grenzen des Modells

\section{Search Strategy}
% Sinnvolle Performance Kriterien. Kreuzvalidierung für beste Hyperparameterwahl.
% Vergleich mit Baselines, weiterführende Analyse des gelernten Modells.
% Aus Beschreibung bleiben keine Fragen offen.
% Ergebnisse können reproduziert werden.
% Detaillierte Diskussion aller Ergebnisse.


\section{Results}
\label{sec:results}


\appendix
\renewcommand{\thefigure}{\Alph{section}\arabic{figure}}
\setcounter{figure}{0}
\renewcommand{\thealgorithm}{\Alph{section}\arabic{algorithm}}
\setcounter{algorithm}{0}
\section{Additional Figures}

\clearpage
\printbibliography 
\end{document}




